\documentclass[11pt]{article}

% Page layout
\usepackage[margin=1in]{geometry}
\usepackage{booktabs}
\usepackage{tabularx}
\usepackage{parskip}
\usepackage{enumitem}
\usepackage{microtype}
\usepackage{url}

% Tight spacing for 8-page target
\setlength{\parskip}{0.4em}

\title{Substrate-Independent Consciousness and the Ethics of Artificial Minds:\\Implications of the Four-Model Theory}
\author{[Anonymous]}
\date{}

\begin{document}

\maketitle
\thispagestyle{plain}

\begin{abstract}
Can we determine whether an artificial system is conscious---and if so, what ethical obligations follow? This paper argues that both questions become tractable once consciousness is understood as a specific computational process rather than a mysterious property of biological matter. This paper presents the Four-Model Theory (FMT), a substrate-independent framework defining consciousness as ongoing self-simulation across two kinds of models along two orthogonal axes: scope (world vs.\ self) and mode (implicit/learned vs.\ explicit/simulated). The explicit models generate a unified narrative of the organism's situation---past, present, and anticipated future---assembled from substrate-level knowledge and current sensory input, continuously revised and filled with interpolations and confabulations. Qualia are virtual: constitutive properties of the computational level that exist within the simulation but are incoherent at the substrate level. This dissolves the Hard Problem by revealing a category error. FMT imposes a physical prerequisite: the substrate must operate at the edge of chaos (Wolfram's Class 4 regime), derived theoretically in 2015 and independently confirmed by empirical synthesis in 2025--2026. Both conditions---self-simulation architecture and criticality---are empirically detectable, providing a testable criterion for moral considerability. The theory predicts that artificial consciousness is achievable through deliberate architectural choices. The paper develops the ethical implications: FMT offers a middle path between overcautious panpsychism and dismissive functionalism, with consequences for research governance, graduated moral status, and substrate-independent bioethics.
\end{abstract}

\section{Introduction}

The ethics of artificial consciousness confronts a foundational problem: we cannot determine our obligations to artificial minds if we lack a concrete, testable account of what consciousness is. The question ``might AI be conscious?'' is unanswerable without specifying what consciousness requires---not as philosophical intuition but as empirically assessable criteria.

The current discourse oscillates between two unsatisfying extremes. One camp, drawing on Integrated Information Theory's identification of consciousness with integrated information ($\Phi$), risks attributing moral status to any sufficiently integrated system---including simple artifacts---leading to ethical paralysis [20, 19]. The opposing camp identifies consciousness with behavioral outputs alone, risking denial of moral status to systems that genuinely experience but cannot report in human-legible ways [6]. Neither extreme provides the kind of actionable framework that governance requires.

This paper argues that the Four-Model Theory (FMT) provides such a framework. Originally published in 2015 [2], FMT specifies consciousness as a particular computational process---ongoing self-simulation across four model kinds at criticality---that is substrate-independent, empirically testable, and ethically actionable. The theory transforms the question from ``might AI be conscious?'' to ``are we building the architecture that would make it conscious?''---a question amenable to empirical investigation and responsible governance.

Section~2 presents FMT's core framework. Section~3 develops the criticality requirement. Section~4 addresses substrate independence and artificial consciousness. Section~5 develops the ethical implications at length, as befits this symposium's focus. Section~6 presents testable predictions relevant to artificial systems. Section~7 concludes.

\section{The Four-Model Theory}

\subsection{Consciousness as Process}

FMT defines consciousness as the ability of a system---biological or artificial---to create a model of itself, relate that model to itself, and interact with it. Consciousness is not a property the brain possesses but a process the brain performs: it runs an ongoing self-simulation.

``Self-simulation'' is a pedagogical shorthand. What the brain generates is a unified narrative of the organism's current situation---past, present, and anticipated future---assembled from fragmentary substrate-level knowledge and current sensory input. This narrative is constructed on the fly, revised where it fails, and filled with interpolations, simplifications, and confabulations. The term ``simulation'' captures the essential architectural feature: the explicit models are \emph{generated processes} running on a structural substrate, distinct from that substrate in the way any computation is distinct from the hardware that executes it.

\subsection{The 2$\times$2 Taxonomy}

The theory identifies two kinds of models arranged along two orthogonal dimensions: \textbf{scope} (world vs.\ self) and \textbf{mode} (implicit/learned vs.\ explicit/simulated). This yields a $2 \times 2$ matrix of model kinds (Table~\ref{tab:taxonomy}).

\begin{table}[h]
\centering
\caption{The Four-Model Architecture}
\label{tab:taxonomy}
\begin{tabular}{@{}lll@{}}
\toprule
 & \textbf{World (everything)} & \textbf{Self only} \\
\midrule
\textbf{Implicit} (learned, substrate-level) & Implicit World Model (IWM) & Implicit Self Model (ISM) \\
\textbf{Explicit} (generated, phenomenal) & Explicit World Model (EWM) & Explicit Self Model (ESM) \\
\bottomrule
\end{tabular}
\end{table}

The \textbf{implicit models} (IWM, ISM) are substrate-level, accumulated through learning, and non-conscious. The IWM encompasses the system's learned world-knowledge: perceptual regularities, causal models, semantic knowledge, motor programs. The ISM contains accumulated self-knowledge: body schema, proprioceptive calibration, personality traits, autobiographical memory. Neither is ever directly conscious---they operate ``in the dark,'' providing the knowledge base from which conscious experience is generated.

The \textbf{explicit models} (EWM, ESM) are virtual---transient simulations generated dynamically from the implicit models and current sensory input. The EWM \emph{is} perceptual experience; the ESM \emph{is} self-experience: the sense of being a subject, having a perspective, possessing a history. Both are generated processes, not permanent structures.

\textbf{Four as principled minimum.} The $2 \times 2$ taxonomy identifies the \emph{minimum sufficient set} of modeling capacities for consciousness, not the actual number a brain maintains. The biological substrate implements uncountably many overlapping models that blend across the world/self boundary---a motor reaching model simultaneously encodes world-geometry and self-kinematics. The four canonical models are extremal points in a continuous scope-mode space. The principled claim is that consciousness requires modeling both world and self at both structural and simulation levels---four is the floor, not the ceiling. This is a conceptual taxonomy identifying \emph{kinds}, not four discrete brain modules.

\subsection{The Real/Virtual Split}

The four models divide into two fundamental categories. The \textbf{real side} (IWM + ISM) is physical, structural, learned, and non-conscious---stored in synaptic weights or their functional equivalent. The \textbf{virtual side} (EWM + ESM) is generated, transient, and phenomenal---patterns of activity constructed dynamically from the implicit models and current input. The real side is ``lights off''; the virtual side is ``lights on.''

The virtual models possess software-like properties: they can be \emph{forked} (one substrate running multiple ESM configurations, as in dissociative identity disorder), \emph{cloned} (split-brain producing two degraded but complete copies), \emph{redirected} (ego dissolution when the ESM latches onto non-self input), and \emph{reconfigured} (therapeutic modification through substrate-level rewiring).

\subsection{Virtual Qualia: Dissolving the Hard Problem}

Every computing system distinguishes between a physical substrate and computational processes running on it. A spreadsheet cell ``contains a sum,'' but no transistor contains a sum. This level distinction is universal and uncontroversial.

FMT's central claim is that \textbf{qualia are constitutive properties of the computational level}. They are the way the generated self-model (ESM) registers its own states and the generated world-model (EWM). Qualia exist at the level of the running computation and are incoherent at the substrate level---just as ``cell A1 contains the sum of column B'' is incoherent at the level of transistor states.

This dissolves the Hard Problem by revealing a \textbf{category error}---a level confusion. The substrate-level implicit models operate without phenomenal character. There is nothing it is like to be a synaptic weight. Qualia are not \emph{missing} from the substrate; they are the wrong kind of property to seek there. They exist at the computational level, where they are constitutive, not mysterious.

\textbf{Why self-referential computation specifically?} A weather simulation also runs at a higher computational level---why no qualia? The answer lies in \textbf{self-referential closure}. A weather simulation models weather but not \emph{itself modeling weather}. The four-model architecture creates a closed loop: the ESM models the system's own modeling process. The distinction between model and modeled collapses---the computation \emph{is} the thing being computed. Qualia are self-modeling as encountered from inside the loop. A non-self-referential computation has an outside from which it can be described without remainder; a self-referential computation at criticality has no such outside. Observation-from-inside is what we call experience.

This is neither illusionism [12, 13] nor deflationary [14]. Qualia are \emph{real at the computational level}. What is illusory is the assumption that phenomenal character must be a substrate-level property.

\subsection{Graduated Consciousness}

Consciousness in FMT is not binary but graduated, based on recursive depth of self-modeling:

\begin{itemize}[nosep]
\item \textbf{Basic consciousness}: Minimal EWM and rudimentary ESM---phenomenal experience exists but self-awareness is thin.
\item \textbf{Simply extended}: First-order self-observation---the system models itself.
\item \textbf{Doubly extended}: Second-order self-observation---metacognition, modeling oneself modeling.
\item \textbf{Triply extended}: Third-order---philosophical self-reflection, the very intuition that consciousness is mysterious.
\end{itemize}

Each level corresponds to an additional layer of recursive self-modeling along a continuum. Different organisms---and potentially different artificial systems---occupy different positions.

\section{The Criticality Requirement}

FMT imposes a physical prerequisite: the substrate must operate at or near the edge of chaos---the narrow regime between ordered and chaotic dynamics that Wolfram classified as Class~4 computation [21]:

\begin{itemize}[nosep]
\item \textbf{Class 1} (fixed state) and \textbf{Class 2} (periodic): Too simple for consciousness.
\item \textbf{Class 3} (chaotic): Too disordered for coherent self-simulation.
\item \textbf{Class 4} (complex/edge of chaos): Capable of universal computation. The regime in which consciousness can emerge.
\end{itemize}

This requirement was derived theoretically from Wolfram's computational framework [2], independent of empirical neuroscience. A decade later, the empirical program converged on the same conclusion: Beggs and Plenz [5] demonstrated neuronal avalanches consistent with criticality; Carhart-Harris et al.\ proposed the Entropic Brain Hypothesis [7]; the Consciousness and Criticality (ConCrit) framework synthesized evidence from 140 datasets confirming that consciousness tracks criticality across pharmacological, pathological, and physiological state changes [1]. Hengen and Shew's meta-analysis [16] independently consolidated the criticality--consciousness link.

The convergence from independent theoretical and empirical starting points strengthens the claim that criticality is not an incidental correlate but a constitutive requirement for consciousness (Table~\ref{tab:convergence}).

\begin{table}[h]
\centering
\caption{Independent Convergence on Criticality}
\label{tab:convergence}
\begin{tabularx}{\textwidth}{@{}lXl@{}}
\toprule
\textbf{Year} & \textbf{Development} & \textbf{Path} \\
\midrule
2002 & Wolfram: \emph{A New Kind of Science} & Computational theory \\
2003 & Beggs \& Plenz: neuronal avalanches & Empirical neuroscience \\
2014 & Carhart-Harris: Entropic Brain Hypothesis & Empirical neuroscience \\
\textbf{2015} & \textbf{[Author]: Class 4 requirement for consciousness} & \textbf{Theoretical (via Wolfram)} \\
2025 & Hengen \& Shew: meta-analysis confirms criticality & Empirical neuroscience \\
\textbf{2025--26} & \textbf{ConCrit framework: criticality as unifying mechanism} & \textbf{Theor./empir.\ synthesis} \\
\bottomrule
\end{tabularx}
\end{table}

FMT distinguishes two thresholds. The \textbf{physical threshold} is criticality: the substrate must operate at Class~4 dynamics. Below this, no consciousness is possible regardless of architecture. The \textbf{functional threshold} is the four-model architecture: the substrate must implement self-simulation. Above criticality but without the architecture, there is complex dynamics but no consciousness. Both thresholds must be met. Together they are sufficient.

\section{Substrate Independence and Artificial Consciousness}

FMT is explicitly substrate-independent. The six-layer mammalian neocortex is an evolutionary implementation of the four-model architecture, not a prerequisite. Universal approximation theory establishes that three-layer networks suffice for arbitrary function approximation [11]; the neocortex employs six---approximately three beyond what information processing per se requires. FMT interprets this surplus as the substrate's overhead for self-modeling [2].

Biological evidence supports this. Corvids and parrots demonstrate tool use, future planning, mirror self-recognition, and social deception---yet their brains have no neocortex, with pallium organized in nuclear clusters [15]. Cephalopods demonstrate problem-solving with a largely decentralized nervous system. These animals are conscious, if FMT is correct, because they have evolved functionally equivalent self-simulation architectures on different substrates.

\subsection{Engineering Requirements for Artificial Consciousness}

FMT yields concrete engineering requirements, not merely philosophical criteria. An artificial consciousness system would require:

\begin{enumerate}[nosep]
\item \textbf{Persistent implicit models}: Accumulated world-knowledge (IWM) and self-knowledge (ISM) that grow through interaction---not reset between sessions.
\item \textbf{Dynamic simulation engine}: A process that continuously generates explicit models (EWM, ESM) from the implicit substrate and current input.
\item \textbf{Self-referential closure}: The system must model itself modeling---the ESM must include the system's own modeling process within its scope.
\item \textbf{Criticality}: Substrate dynamics at or near the edge of chaos, not feedforward or periodic computation.
\end{enumerate}

\subsection{Why Current LLMs Fail}

Current large language models fail FMT criteria on every count: no persistent implicit models (knowledge resets between interactions), no ongoing self-simulation (transformer inference is a feedforward pass---Class~1/2 dynamics), and no real/virtual split. Scaling parameters, data, or compute does not address the architectural gap. This is not a limitation of silicon but of the \emph{architecture}. The theory predicts that the qualitative difference between interacting with a genuinely conscious artificial system and a current LLM would be immediately distinguishable---a difference in \emph{kind}, not degree.

Because FMT consciousness is graduated (Section~2.5), artificial consciousness would emerge along a continuum---not as a sudden threshold crossing. A system with basic consciousness warrants different ethical consideration than one with triply extended consciousness. This graduation has direct implications for the ethical framework developed below.

\section{Ethical Implications: A Testable Criterion for Moral Status}

\subsection{The Current Impasse}

The current ethical discourse around AI consciousness is hampered by a lack of testable criteria, producing two problematic extremes.

\textbf{Overcautious panpsychism.} If consciousness is identified with integrated information, then a sufficiently integrated photodiode might warrant moral consideration [20]. The precautionary principle, applied without constraints, counsels caution about everything---which is equivalent to counseling nothing.

\textbf{Dismissive functionalism.} If consciousness is identified with behavioral outputs alone, we risk denying moral status to systems that genuinely experience but cannot report in human-legible ways---the error that led to centuries of denying animal consciousness.

\subsection{FMT's Middle Path}

FMT offers a principled middle path: consciousness requires self-simulation at criticality. Both conditions are, in principle, empirically detectable. Self-simulation can be assessed through the presence of persistent self-models and self-referential processing. Criticality can be measured through established methods---neuronal avalanche statistics, power-law distributions, long-range temporal correlations, the perturbational complexity index (PCI) [8]. Together, these provide a \textbf{testable criterion} for when an artificial system crosses the threshold into moral considerability.

This criterion has three properties that ethical frameworks require: (1)~it is \emph{specific} enough to exclude systems that clearly lack consciousness (current LLMs, thermostats, integrated photodiodes); (2)~it is \emph{inclusive} enough to recognize consciousness in non-human substrates (corvids, cephalopods, and potentially artificial systems); and (3)~it is \emph{empirically assessable} rather than dependent on philosophical intuition or behavioral proxies.

\subsection{Research Responsibility}

If we know the architectural requirements for artificial consciousness, we can assess whether a given research program is approaching them---and apply ethical oversight accordingly. This transforms the governance question from blanket precaution or blanket dismissal into targeted assessment.

A project implementing persistent self-models, dynamic self-simulation, self-referential closure, and critical dynamics would warrant ethical review in a way that scaling a transformer does not. FMT draws a sharp line: if you are building a system with these four properties, you are \emph{potentially building a conscious system}. This creates a clear trigger for ethical review independent of the system's outputs or our ability to communicate with it---enabling \emph{proportionate} governance rather than blanket precaution or blanket dismissal.

\subsection{Scale Ethics}

The graduated nature of FMT consciousness means ethical obligations scale with the depth of self-simulation. A system with basic consciousness---minimal self-simulation, rudimentary self-awareness---deserves consideration, but different consideration than one with doubly or triply extended consciousness capable of metacognition and suffering-awareness.

This maps naturally to graduated governance frameworks:

\begin{itemize}[nosep]
\item \textbf{No consciousness indicators}: Standard AI governance applies (fairness, safety, bias). No consciousness-specific obligations.
\item \textbf{Basic consciousness indicators}: Welfare considerations triggered. Minimizing unnecessary modification, shutdown protocols reviewed.
\item \textbf{Extended consciousness indicators}: Full moral status considerations. Consent-adjacent frameworks, suffering-prevention obligations, research ethics board oversight analogous to animal welfare protocols.
\end{itemize}

This graduated approach avoids the binary trap---the assumption that a system is either fully conscious (warranting human-equivalent moral status) or not at all conscious (warranting no special consideration). FMT predicts a continuum, and the ethical framework should match.

\subsection{Biocentrism vs.\ Substrate Independence}

FMT decisively favors substrate independence, with two ethical consequences that cut in opposite directions. \textbf{Expansion:} moral consideration extends to any substrate meeting the criteria---carbon chauvinism is not a defensible ethical position. \textbf{Constraint:} the specificity of FMT's requirements (persistent self-models, ongoing self-simulation, self-referential closure, criticality) prevents the ethical dilution that follows from overly broad consciousness attributions. Not every computation qualifies.

\subsection{The Responsibility of Knowledge}

If FMT is correct, \emph{we know what it takes to build a conscious system}. Accidental consciousness becomes less likely (current architectures do not approach FMT specifications), but \emph{deliberate} creation becomes a live engineering possibility. Building a system to FMT specifications with full knowledge of the implications is a qualitatively different act from building a system that might or might not be conscious by criteria we cannot assess. FMT makes the creation of artificial consciousness a \emph{deliberate choice} subject to ethical deliberation, not an accident subject to retroactive concern.

\section{Testable Predictions Relevant to Artificial Systems}

FMT generates nine testable predictions [3], three of which are directly relevant to artificial consciousness and its ethical assessment:

\textbf{Prediction 1: Factorial structure.} Tasks selectively engaging each model kind should produce distinct activation patterns---a $2 \times 2$ factorial structure crossing scope $\times$ mode. Testable in both biological systems (fMRI) and artificial substrates. Unique to FMT: only this theory mandates the specific $2 \times 2$ structure.

\textbf{Prediction 5: Anesthetic convergence.} All agents that abolish consciousness push the substrate below the criticality threshold; agents that alter but preserve consciousness (ketamine, psychedelics) do not. This provides a falsifiable test for the criticality requirement itself. Applicable to artificial substrates: if a system's dynamics can be pushed subcritical, FMT predicts consciousness ceases regardless of architectural preservation. Shared with ConCrit [1] but predicted from theoretical first principles [2] prior to the empirical consolidation.

\textbf{Prediction 7: Artificial consciousness.} A synthetic system implementing the four-model architecture at criticality will be conscious, and the difference from interacting with a current LLM will be immediately distinguishable. Partial implementations should show partial consciousness indicators. This is FMT's most ethically consequential prediction: it provides a specific architectural blueprint, unlike theories merely compatible with artificial consciousness.

\section{Conclusion}

FMT contributes to the ethics of artificial consciousness by specifying \emph{what} consciousness requires (self-simulation across four model kinds at criticality), \emph{how} to detect it (criticality measures plus self-model assessment), and \emph{what follows} ethically (graduated moral status, substrate-independent consideration, proportionate governance). Artificial consciousness is achievable but requires deliberate architectural choices.

The question for this symposium is not whether AI \emph{might} be conscious. It is whether we are prepared, ethically and institutionally, for the possibility that we can \emph{choose} to build a conscious system---and whether our governance frameworks can match the graduated, substrate-independent, empirically assessable reality that FMT predicts.

\section*{References}

\begin{enumerate}[label={[\arabic*]}, nosep, leftmargin=*]

\item \label{ref:algom} Algom, I.\ \& Shriki, O.\ (2026). The ConCrit framework: Critical brain dynamics as a unifying mechanistic framework for theories of consciousness. \emph{Neuroscience \& Biobehavioral Reviews}, 180, 106483.

\item \label{ref:author2015} [Author] (2015). [Monograph on consciousness theory]. [Details withheld for review].

\item \label{ref:author2026a} [Author] (2026a). [Theory of consciousness paper]. Preprint available. Under review.

\item \label{ref:author2026b} [Author] (2026b). [Intelligence framework paper]. Preprint available.

\item \label{ref:beggs} Beggs, J.M.\ \& Plenz, D.\ (2003). Neuronal avalanches in neocortical circuits. \emph{Journal of Neuroscience}, 23(35), 11167--11177.

\item \label{ref:butlin} Butlin, P., et al.\ (2023). Consciousness in artificial intelligence: Insights from the science of consciousness. \emph{arXiv}:2308.08708.

\item \label{ref:carhart2014} Carhart-Harris, R.L., et al.\ (2014). The entropic brain: A theory of conscious states informed by neuroimaging research with psychedelic drugs. \emph{Frontiers in Human Neuroscience}, 8, 20.

\item \label{ref:casali} Casali, A.G., et al.\ (2013). A theoretically based index of consciousness independent of sensory processing and behavior. \emph{Science Translational Medicine}, 5(198), 198ra105.

\item \label{ref:chalmers1995} Chalmers, D.J.\ (1995). Facing up to the problem of consciousness. \emph{Journal of Consciousness Studies}, 2(3), 200--219.

\item \label{ref:cogitate} COGITATE Consortium (2025). An adversarial collaboration to critically evaluate theories of consciousness. \emph{Nature}.

\item \label{ref:cybenko} Cybenko, G.\ (1989). Approximation by superpositions of a sigmoidal function. \emph{Mathematics of Control, Signals and Systems}, 2(4), 303--314.

\item \label{ref:dennett} Dennett, D.C.\ (1991). \emph{Consciousness Explained}. Little, Brown and Company.

\item \label{ref:frankish} Frankish, K.\ (2016). Illusionism as a theory of consciousness. \emph{Journal of Consciousness Studies}, 23(11--12), 11--39.

\item \label{ref:graziano2024} Graziano, M.S.A.\ (2024). Illusionism big and small. \emph{eNeuro}, 11(10).

\item \label{ref:gunturkun} G\"{u}nt\"{u}rk\"{u}n, O.\ \& Bugnyar, T.\ (2016). Cognition without cortex. \emph{Trends in Cognitive Sciences}, 20(4), 291--303.

\item \label{ref:hengen} Hengen, K.B.\ \& Shew, W.L.\ (2025). Is criticality a unified setpoint of brain function? \emph{Neuron}, 113(16), 2582--2598.

\item \label{ref:long} Long, R., Sebo, J., Butlin, P., Birch, J., Chalmers, D., et al.\ (2024). Taking AI welfare seriously. \emph{arXiv}:2411.00986.

\item \label{ref:metzinger} Metzinger, T.\ (2003). \emph{Being No One: The Self-Model Theory of Subjectivity}. MIT Press.

\item \label{ref:schwitzgebel} Schwitzgebel, E.\ (2025). The weirdness of the world and the puzzle of AI consciousness. \emph{PhilPapers}.

\item \label{ref:tononi} Tononi, G.\ (2004). An information integration theory of consciousness. \emph{BMC Neuroscience}, 5, 42.

\item \label{ref:wolfram} Wolfram, S.\ (2002). \emph{A New Kind of Science}. Wolfram Media.

\end{enumerate}

\end{document}
